\documentclass{article}
\usepackage[margin=1in]{geometry}
\usepackage{amsmath,amssymb}
\usepackage[most]{tcolorbox}

\newcommand{\BookName}{Pattern Recognition and Machine Learning}
\newcommand{\BookAuthor}{Christopher M. Bishop}
\newcommand{\ChapterName}{Introduction}
\newcommand{\ExerciseID}{1.7}

\title{Chapter: \ChapterName}
\author{Ahonakpon Guy Baguidi}
\date{}

\begin{document}

\maketitle

\begin{center}
  \large\textbf{Exercise \ExerciseID}

  \vspace{0.5em}

  \normalsize\BookName\ - \BookAuthor
\end{center}

\begin{tcolorbox}[breakable,colback=gray!5,colframe=gray!60,title=Solution]
In this exercise, we prove the normalization condition (1.48)
for the univariate Gaussian distribution:

\begin{align*}
  \int_{-\infty}^{\infty} N(x|\mu, \sigma^2) dx = 1
\end{align*}

where $N(x|\mu, \sigma^2)$ is the probability density function of a univariate Gaussian distribution with mean $\mu$ and variance $\sigma^2$:

\begin{align*}
  N(x|\mu, \sigma^2) = \frac{1}{\sqrt{2\pi\sigma^2}} \exp\left( -\frac{(x-\mu)^2}{2\sigma^2} \right)
\end{align*}

To do this consider, the integral:

\begin{align*}
  I &= \int_{-\infty}^{\infty}  \exp\left( -\frac{(x-\mu)^2}{2\sigma^2} \right) dx
\end{align*}

which we can evaluate by first writing its square in the form:

\begin{align*}
  I^2 &= \int_{-\infty}^{\infty} \int_{-\infty}^{\infty}  \exp\left( -\frac{1}{2\sigma^2} x^2 -\frac{1}{2\sigma^2} y^2 \right) dx dy
\end{align*}

Now make the transformation from cartesian coordinates $(x, y)$ to polar coordinates $(r, \theta)$ and then 
substitute $u = r^2$ to obtain:

Let's show that, by performing the integrals over $\theta$ and $u$, and then taking the square root of both sides,
we obtain 

\begin{align*}
  I &= {2\pi\sigma^2}^{1/2}
\end{align*}

Let's start by performing the change of variables from cartesian coordinates $(x, y)$ to polar coordinates $(r, \theta)$:

\begin{align*}
  I^2 &= \int_{0}^{2\pi} \int_{0}^{\infty}  \exp\left( -\frac{1}{2\sigma^2} r^2 \right) r dr d\theta
\end{align*}

Next, we perform the substitution $u = r^2$, which implies that $du = 2r dr$ or equivalently $r dr = \frac{1}{2} du$:

\begin{align*}
  I^2 &= \int_{0}^{2\pi} \int_{0}^{\infty}  \exp\left( -\frac{1}{2\sigma^2} u \right) \frac{1}{2} du d\theta \\
  &= \frac{1}{2} \int_{0}^{2\pi} d\theta \int_{0}^{\infty}  \exp\left( -\frac{1}{2\sigma^2} u \right) du \\
  &= \pi \int_{0}^{\infty}  \exp\left( -\frac{1}{2\sigma^2} u \right) du \quad \text{since } \int_{0}^{2\pi} d\theta = 2\pi \text{ and } \frac{1}{2} \cdot 2\pi = \pi
\end{align*}

Now we can evaluate the remaining integral over $u$:

\begin{align*}
  \int_{0}^{\infty}  \exp\left( -\frac{1}{2\sigma^2} u \right) du &= \left[ -2\sigma^2 \exp\left( -\frac{1}{2\sigma^2} u \right) \right]_{0}^{\infty} \\
  &= 2\sigma^2
\end{align*}

Substituting this result back into the expression for $I^2$, we have:

\begin{align*}
  I^2 &= \pi \cdot 2\sigma^2 \\
  &= 2\pi\sigma^2
\end{align*}

Finally, by taking the square root of both sides, we obtain:

\begin{align*}
  I &= \sqrt{2\pi\sigma^2} \\
  &= {2\pi\sigma^2}^{1/2}
\end{align*}

Finally by using this result, we show that the normalization condition holds:

\begin{align*}
  \int_{-\infty}^{\infty} N(x|\mu, \sigma^2) dx &= \int_{-\infty}^{\infty} \frac{1}{\sqrt{2\pi\sigma^2}} \exp\left( -\frac{(x-\mu)^2}{2\sigma^2} \right) dx \\
  &= \frac{1}{\sqrt{2\pi\sigma^2}} \int_{-\infty}^{\infty}  \exp\left( -\frac{(x-\mu)^2}{2\sigma^2} \right) dx \\
  &= \frac{1}{\sqrt{2\pi\sigma^2}} I \\
  &= \frac{1}{\sqrt{2\pi\sigma^2}} {2\pi\sigma^2}^{1/2} \\
  &= 1
\end{align*}

\par\medskip
\paragraph{Personal note:} During the process of solving this exercise, I learned how to
use the change of variables in multiple integrals (because of the integral operator properties) 
and how to evaluate the Gaussian integral by using polar coordinates and substitution.

\end{tcolorbox}

\end{document}
